\scanpage{65}
\mspage{58}

$W$ admet une structure de \unsure{Qu.f.} \inserted{$(C,\mathrm{TC},F,\mathrm{TF})$} qui soit close et exacte à gauche -- alors les conditions de la définition 1 sont satisfaites. \deleted{\unsure{La p}}

D'autre part, notons que la prop. 1 page \unsure{45|55}, et ses corollaires, y inclus la prop 3 (p. 50) et son corollaire, sont \deleted{\unsure{satisfait}} valables dès lors que les quillénisations des $(M(I),W(I))$ envisagées sont définies ``uniformément à gauche'' par une quillénisation uniforme à gauche (pas \unsure{néct|nécess.} ``exacte à gauche'') de $(M,W)$ sur $\mathrm{Diag}$.

La quillénisation uniforme gauche nous servira à établir l'existence des images directes homotopiques $u_{\unsure{*}}$, pour certains $u$. La quillénisation uniforme droite, \unsure{dualement}, \unsure{quand|que} nous \unsure{considérons|considérerons} \unsure{principalt|principalement|présentement} l'\unsure{existce|existence} des $u_{\unsure{!|*}}$ \inserted{homotopiques}, pour certains $u:I'\to I$ dans $\mathrm{Diag}$.

Je vais résumer ce que j'ai pu établir jusqu'à présent dans ce sens, dans \unsure{la|le|les}
% archive sheet number: 64
